\documentclass[11pt,a4paper]{article}

\usepackage[T1]{fontenc}
\usepackage[utf8]{inputenc}
\usepackage[brazil]{babel}
\usepackage{lmodern}
\usepackage{microtype}
\usepackage[margin=2cm]{geometry}
\usepackage{hyperref}
\usepackage{titlesec}
\usepackage{xcolor}

\hypersetup{
  colorlinks=true,
  linkcolor=black,
  urlcolor=black,
  citecolor=black,
  pdfborder={0 0 0}
}

\definecolor{sectioncolor}{HTML}{111111}
\definecolor{accentcolor}{HTML}{1A1A1A}

\pagestyle{empty}
\setlength{\parindent}{0pt}
\setlength{\parskip}{5pt}
\linespread{1.05}
\selectfont

\titleformat{\section}
  {\bfseries\large\color{sectioncolor}}
  {}{0pt}{}[\vspace{-2pt}{\color{sectioncolor}\titlerule[0.8pt]}]
\titlespacing*{\section}{0pt}{12pt}{8pt}

\newcommand{\cvrole}[4]{%
  \vspace{4pt}%
  {\bfseries\color{accentcolor}#1}\hfill{\itshape #2}\\[-2pt]%
  {\itshape #3}\hfill #4\\[6pt]%
}

\begin{document}

\begin{center}
  {\fontsize{20}{24}\selectfont\bfseries Kamille Hartmann Maccagnan}\\[8pt]
  {\large Atendimento Corporativo \textbar{} Suporte de TI \textbar{} Indicadores e Processos}\\[10pt]
  Campo Comprido, Curitiba-PR \enspace\textbullet\enspace (41) 9 9928-4424 \enspace\textbullet\enspace
  \href{mailto:kamillehartmannm@gmail.com}{kamillehartmannm@gmail.com}\\[3pt]
  \href{https://linkedin.com/in/kamille-hartmann-maccagnan/}{linkedin.com/in/kamille-hartmann-maccagnan/}
\end{center}

\vspace{6pt}

\section{Resumo Profissional}

Profissional em formação na área de tecnologia, com experiência em suporte a usuários, atendimento corporativo e apoio operacional em ambientes de alto volume de demandas. Atua na orientação e capacitação de equipes, padronização de processos, monitoramento de indicadores de desempenho e comunicação clara com usuários e áreas técnicas por e-mail, registros formais e interações presenciais. Vivência com hardware e software na perspectiva do usuário, documentação de procedimentos, gestão de chamados no ServiceNow e análise de indicadores com Excel, Google Sheets e Power BI. Perfil comunicativo, empático, analítico e orientado a resultados, com proatividade para priorizar entregas e promover melhoria contínua no atendimento.

\section{Experiência Profissional}

\cvrole{Anjuss}{Analista de Suporte de TI}{07/2025 -- Atual}

Atuo no suporte a usuários e sistemas da rede de lojas, com atendimento corporativo, organização e priorização de demandas em ambiente dinâmico e de alto volume de solicitações. Lidero treinamentos e orientação de novos funcionários, assegurando padronização, alinhamento e qualidade no atendimento e na utilização das ferramentas. Elaboro e atualizo documentações e POPs para lojas e equipe de T.I., garantindo clareza nos fluxos e rastreabilidade das rotinas. Realizo testes manuais do sistema antes da liberação para as lojas, reportando ocorrências com comunicação objetiva à equipe técnica. Desenvolvo dashboards em Power BI para monitoramento de indicadores de desempenho do atendimento e dos processos, propondo ações de melhoria com base na análise dos resultados.

\cvrole{HORSE do Brasil}{Estagiária de TI (IS/IT LATAM)}{07/2024 -- 07/2025}

Atuei no suporte e acompanhamento de demandas em ambiente multinacional, utilizando o ServiceNow para registro, acompanhamento e reporte de incidentes, com foco em agilidade, eficiência e satisfação dos usuários internos. Apoiei a equipe e a coordenação de atividades com documentação de processos, verificação de informações e comunicação formal de ocorrências para resolução. Elaborei relatórios e dashboards em Power BI para monitoramento de indicadores operacionais, identificando padrões e oportunidades de melhoria no fluxo de atendimento. Atuei como ponte entre usuários, áreas técnicas e stakeholders, participando de reuniões multilíngues e assegurando alinhamento entre as partes envolvidas.

\cvrole{Restaurante Familiar}{Analista de Dados e Suporte Técnico}{01/2023 -- 07/2024}

Atuei no suporte técnico a usuários e na validação de informações operacionais, com atenção aos detalhes e comunicação clara sobre problemas e necessidades de correção. Estruturei bases de dados e dashboards em Power BI e Excel para acompanhamento de indicadores de desempenho, apoiando a identificação de desvios e a proposição de melhorias nos processos. Interagi com áreas operacionais e administrativas para garantir alinhamento nas demandas, padronização das rotinas e eficiência no relacionamento interno com os usuários do negócio.

\section{Competências Técnicas}

\textbf{Atendimento e Suporte:} Suporte a usuários, atendimento corporativo, hardware e software na perspectiva do usuário, gestão e priorização de demandas, comunicação formal por e-mail e registros, reporte de ocorrências, orientação e capacitação de equipes, foco em satisfação e qualidade.

\textbf{Processos e Indicadores:} Monitoramento de KPIs, análise de indicadores de desempenho, padronização de fluxos, documentação e POPs, melhoria contínua, capacidade analítica, identificação de padrões e proposição de soluções.

\textbf{Ferramentas:} ServiceNow, Excel avançado, Google Sheets, Power BI, Pacote Office, Trello, Monday.

\textbf{Colaboração:} Trabalho em equipe, empatia, comunicação interpessoal, interface entre usuários e áreas técnicas, apoio à coordenação, postura proativa e orientada a resultados.

\textbf{Idiomas:} Inglês Avançado \enspace\textbullet\enspace Espanhol Básico \enspace\textbullet\enspace Coreano Básico

\section{Projetos e Conquistas}

1º lugar no Transformation Day Renault 2023 com solução orientada a melhoria de processos e atendimento. Participação no Hackathon Bosch 2023 com foco em resolução de problemas. Desenvolvimento de documentações, POPs e treinamentos para padronização do atendimento em rede de lojas. Implementação de dashboards para monitoramento de indicadores operacionais e apoio à tomada de decisão.

\section{Formação Acadêmica}

\textbf{PUCPR} \hfill Curitiba, PR\\
Graduação em Gestão da Tecnologia da Informação \hfill Previsão de conclusão: 06/2026\\[6pt]

\textbf{Escola Conquer} \hfill Curitiba, PR\\
Pós-graduação em Liderança e Gestão em Tecnologia \hfill Conclusão: 2024\\[6pt]

\textbf{Universidade Tuiuti do Paraná} \hfill Curitiba, PR\\
Graduação em Medicina Veterinária \hfill 06/2019\\[6pt]

\textbf{Qualittas} \hfill Curitiba, PR\\
Pós-graduação em Clínica Médica e Cirúrgica de Animais Exóticos e Pets Não Convencionais \hfill 12/2022

\end{document}
